%==============================================================================
%  ACADEMIC SLIDES TEMPLATE  (beamer talk + appendix slides)
%------------------------------------------------------------------------------
%  A minimal 16:9 talk to accompany paper-template.tex. Sans-serif throughout
%  (text and math in Source Sans Pro, the paper's sans), with only the
%  packages the slides use. The body is a short example talk followed by
%  appendix slides numbered apart from the main talk. Replace the metadata and
%  the example content with your own.
%==============================================================================
%  COMPILE (pdflatex only: the sans-serif math needs Type 1 fonts):
%      latexmk -pdf filename
%  or run pdflatex twice (the second pass fills in the slide totals).
%
%  HANDOUT: add "handout" to the class options for one page per frame with
%  every overlay shown (for printing or posting).
%==============================================================================

% t: top-align frame content so it does not shift as overlays appear.
\documentclass[
    aspectratio=169,11pt,t,
    %handout,
]{beamer}

%%%% Theme --------------------------------------------------------------------
%
%  Moloch is the maintained fork of Metropolis. Its title layouts are replaced
%  in "Configuration" so every title sits over the same rule.

\usetheme{moloch}

%%%% Font face packages -------------------------------------------------------
%
%  Source Sans Pro for text and math alike. mathastext takes the math letters,
%  digits, and Greek (via the LGR encoding) from the text font;
%  professionalfonts stops beamer from swapping in its own math fonts first.

\usepackage[semibold]{sourcesanspro}
\usepackage[LGR,T1]{fontenc}
\usefonttheme{professionalfonts}
\usepackage[italic,LGRgreek]{mathastext}

%%%% Packages -----------------------------------------------------------------
%
%  Beamer already loads graphicx, xcolor, hyperref, amsmath, and TikZ (via
%  the theme). Add only what your slides use.

\usepackage{booktabs}

%%%% Configuration ------------------------------------------------------------

% Palette. EXACTLY ONE block must be active (uncommented). Each sets the five
% colors the deck uses: paper (background), ink (text, titles), accent (title
% rules, bullets, block shading), highlight (\alert), and closing (background
% of the closing "Thank you" slide). The commented blocks are alternatives set
% aside for later, in order of preference; slides-palettes.tex shows them all.
% Reuse the hex codes in R / Stata / Python figures so they match the slides.

% 0. Dark rules, vermilion alerts   <-- DEFAULT
\definecolor{paper}{HTML}{FAFAFA}
\definecolor{ink}{HTML}{23373B}
\colorlet{accent}{ink}
\definecolor{highlight}{HTML}{D55E00} % Okabe--Ito vermilion
\definecolor{closing}{HTML}{333333}   % charcoal

% 1. Paul Tol high-contrast
%\definecolor{paper}{HTML}{FAFAFA}
%\definecolor{ink}{HTML}{1F1F1F}
%\definecolor{accent}{HTML}{004488}
%\definecolor{highlight}{HTML}{BB5566}
%\colorlet{closing}{ink}

% 2. Monochrome + red
%\definecolor{paper}{HTML}{FFFFFF}
%\definecolor{ink}{HTML}{1A1A1A}
%\definecolor{accent}{HTML}{5C5C5C}
%\definecolor{highlight}{HTML}{C62828}
%\colorlet{closing}{ink}

% 3. Nord
%\definecolor{paper}{HTML}{FAFAFA}
%\definecolor{ink}{HTML}{2E3440}
%\definecolor{accent}{HTML}{5E81AC}
%\definecolor{highlight}{HTML}{BF616A}
%\colorlet{closing}{ink}

% 4. IBM Carbon
%\definecolor{paper}{HTML}{FFFFFF}
%\definecolor{ink}{HTML}{161616}
%\definecolor{accent}{HTML}{0F62FE}
%\definecolor{highlight}{HTML}{DA1E28}
%\colorlet{closing}{ink}

% 5. Solarized light
%\definecolor{paper}{HTML}{FDF6E3}
%\definecolor{ink}{HTML}{073642}
%\definecolor{accent}{HTML}{268BD2}
%\definecolor{highlight}{HTML}{DC322F}
%\colorlet{closing}{ink}

% 6. Okabe--Ito: blue rules, vermilion alerts
%\definecolor{paper}{HTML}{FAFAFA}
%\definecolor{ink}{HTML}{23373B}
%\definecolor{accent}{HTML}{0072B2}
%\definecolor{highlight}{HTML}{D55E00}
%\definecolor{closing}{HTML}{333333}

% Light frame titles on the page background. The closing (standout) slide
% sets paper-colored text on the closing color.
\setbeamercolor{normal text}{fg=ink,bg=paper}
\setbeamercolor{palette primary}{fg=paper,bg=closing}
\setbeamercolor{structure}{fg=accent}
\setbeamercolor{frametitle}{fg=ink,bg=}
\setbeamercolor{block title}{fg=ink,bg=accent!15}
\setbeamercolor{block body}{bg=accent!5}
% Alerts are also bold, because vermilion is a little light for
% normal-weight text on a white slide.
\setbeamercolor{alerted text}{fg=highlight}
\setbeamerfont{alerted text}{series=\bfseries}

% The deck's one rule: thin, accent-colored, the width of the text block, set
% a fixed 6pt below the baseline of the title it underlines (whatever its
% letters) -- on the title page, every frame, and every section divider.
% Frame titles start at the text margin (beamer's default 1cm), so titles,
% rules, and body text all share one left edge.
\newcommand{\titlerule}{%
  \par\kern-\prevdepth\kern6pt\nointerlineskip
  \noindent{\color{accent}\rule{\textwidth}{0.4pt}}\par}
\setbeamertemplate{title}{\raggedright\inserttitle\par}
\setbeamertemplate{title separator}{\titlerule\vspace*{0.8em}}
\setbeamertemplate{frametitle}{%
  \begin{beamercolorbox}[wd=\paperwidth,leftskip=1cm,rightskip=1cm plus 1fil]{frametitle}%
    \rule{0pt}{4ex}\insertframetitle\par
    {\usebeamerfont{framesubtitle}\insertframesubtitle\par}%
    \titlerule
  \end{beamercolorbox}%
}
\setbeamertemplate{section page}{%
  {\usebeamerfont{section title}\insertsectionhead\par}%
  \titlerule
}

% Footer shows "n / N", where N counts the main talk only; the appendix block
% switches to "Appendix k / K".
\setbeamertemplate{page number in head/foot}[appendixframenumber]

% Speaker notes (\note{...}) are left out of the PDF by default. Uncomment to
% put them on a second screen for a presenter tool such as pdfpc.
%\setbeameroption{show notes on second screen=right}

%%%% Paths --------------------------------------------------------------------
%
%  Folder searched by \includegraphics{...}. Relative to this file.
\graphicspath{{figures/}}

%%%% Metadata -----------------------------------------------------------------

\title{Commitment and present bias: Experimental evidence}
\author{First M.~Last}
\institute{Department of Economics, Example University}
\date{Example Seminar Series\\30 June 2026}


\begin{document}

%%%%%%%%%%%%%%%%%%%%%%%%%
%%%% BEGIN MAIN TALK %%%%
%%%%%%%%%%%%%%%%%%%%%%%%%

% The title page is unnumbered; each \section adds an unnumbered divider.
\maketitle

\section{Introduction}

\begin{frame}{Motivation}
  % [<+->] reveals one item per click.
  \begin{itemize}[<+->]
    \item People plan to act patiently, then take the immediate reward when
    the moment arrives (Doe 2020)
    \item \alert{Present bias:} an extra discount on everything that is not
    ``now'' (Lee, Kim, and Park 2021)
    \item People who know they are present-biased seek commitment
  \end{itemize}
\end{frame}

\begin{frame}{This paper}
  \begin{columns}[T]
    \begin{column}{0.5\textwidth}
      \begin{itemize}
        \item Laboratory experiment: intertemporal allocation decisions under
        one of two conditions
        \item Estimate a quasi-hyperbolic discounting model
        \item \alert{Finding:} the treatment lowers present bias and raises
        demand for commitment
      \end{itemize}
    \end{column}
    \begin{column}{0.45\textwidth}
      \includegraphics[width=\linewidth]{example-image}
    \end{column}
  \end{columns}
\end{frame}

\section{Model}

\begin{frame}{Quasi-hyperbolic discounting}
  A decision maker values a stream of consumption as
  \[
    U_t = u(x_t) + \beta \sum_{\tau=1}^{T} \delta^\tau u(x_{t+\tau})
  \]
  \begin{itemize}
    \item $\beta = 1$: exponential discounting
    \item $\beta < 1$: present bias
  \end{itemize}
\end{frame}

\begin{frame}{Hypotheses}
  \begin{block}{Hypothesis 1}
    The treatment lowers present bias:
    $\beta_{\text{treat}} < \beta_{\text{control}}$
  \end{block}
  \pause
  \begin{block}{Hypothesis 2}
    The treatment raises take-up of the commitment device
  \end{block}
\end{frame}

\begin{frame}{Estimation}
  \begin{enumerate}
    \item Estimate $\beta$ and $\delta$ for each participant by maximum
    likelihood
    \item Compare $\beta$ across conditions, with standard errors that
    account for the two-stage estimator
  \end{enumerate}
\end{frame}

\section{Results}

% label= makes the frame a target for the appendix slide's return button.
% A generated table drops in the same way: \input{tables/table_regression.tex}
\begin{frame}[label=results]{Main results}
  \centering
  \begin{tabular}{lcc}
    \toprule
                     & Model 1 & Model 2 \\
    \midrule
    Treatment effect & 0.281   & 0.895   \\
                     & (0.044) & (0.052) \\
    \midrule
    Observations     & 897     & 897     \\
    \bottomrule
  \end{tabular}

  \medskip
  {\footnotesize Standard errors in parentheses.}

  \vfill
  \hfill\hyperlink{robustness}{\beamergotobutton{Robustness}}
\end{frame}

% Build a figure across clicks: export one PDF per step and show each in turn.
\begin{frame}{Commitment demand}
  \centering
  \includegraphics<1>[height=0.75\textheight]{example-image-a}%
  \includegraphics<2>[height=0.75\textheight]{example-image-b}
  \note{Walk through the control group first; the treatment group appears on
  the next click.}
\end{frame}

\section{Conclusion}

\begin{frame}{Takeaways}
  \begin{itemize}
    \item A simple change in how future rewards are presented lowers present
    bias
    \item It also raises demand for commitment
    \item Caveat: a single laboratory sample; field replication is next
  \end{itemize}
\end{frame}

% Standout frames are unnumbered.
\begin{frame}[standout]
  Thank you

  \medskip
  {\normalsize\mdseries
  Paper: \href{https://example.com/paper.pdf}{example.com/paper.pdf}\\
  \href{mailto:you@example.edu}{you@example.edu}\par}
\end{frame}


%%%%%%%%%%%%%%%%%%%%%%%%
%%%% BEGIN APPENDIX %%%%
%%%%%%%%%%%%%%%%%%%%%%%%
%
%  Appendix slides answer anticipated questions. They are left out of the
%  main slide count and numbered "Appendix k / K". Link to one from the main
%  talk with \hyperlink{<frame label>}{\beamergotobutton{...}} and back with
%  \beamerreturnbutton. Delete this whole block if the talk has no appendix.

% A plain name for the appendix's PDF bookmark; beamer's default goes through
% \translate, which hyperref warns about.
\renewcommand{\appendixname}{Appendix}
\appendix
\setbeamertemplate{page number in head/foot}{%
  Appendix \insertframenumberinappendix\,/\,\insertappendixframenumber}

\begin{frame}[label=robustness]{Robustness: excluding inattentive participants}
  \centering
  \begin{tabular}{lcc}
    \toprule
                     & Model 1 & Model 2 \\
    \midrule
    Treatment effect & 0.264   & 0.871   \\
                     & (0.049) & (0.058) \\
    \midrule
    Observations     & 812     & 812     \\
    \bottomrule
  \end{tabular}

  \medskip
  {\footnotesize Standard errors in parentheses.}

  \vfill
  \hfill\hyperlink{results}{\beamerreturnbutton{Main results}}
\end{frame}

\end{document}
